\section{An upper bound for $C_6$}\label{sec:C6-ub}
In this section we prove Theorem~\ref{thm:C6-ub}, which is restated next.

% \medskip
% \noindent{\textbf{Theorem~\ref{thm:C6-ub}}~
% \emph{\CsixUB}
% }
\begin{theorem}\label{thm:C6-ub}
% \CsixUB
There exists a one-sided error algorithm for testing $C_6$-freeness in graphs of constant arboricity  whose query complexity and running time are $\tilde{O}(n^{1/2}\cdot {\rm poly}(1/\eps))$.
%\commentt{exponent?}
\end{theorem}

\iffalse
\begin{theorem}\label{thm:C6-ub}
There exists a one-sided error algorithm for testing $C_6$-freeness in graphs of constant arboricity  whose query complexity is $\tilde{O}(n^{1/2}/\eps^{7})$. \commentt{DR: check exponent - maybe write as $\tilde{O}(n^{1/2})\cdot \cdot \poly(1/\eps)$?}
\end{theorem}
\fi

The algorithm referred to in Theorem~\ref{thm:C6-ub} is presented next.


\medskip\noindent
\textsf{Test-$C_6$-freeness}($n,\epsilon,\alpha$)\label{alg:C6}
\begin{enumerate}
\item Let $\theta_0 = 4 \alpha/\eps$, and $\theta_1 = c_1\cdot n^{1/2}\log^2 n/\eps^2$ 
(for a sufficiently large constant $c_1$).
%\commentt{DR: set $c'$ and $\theta_1$, looks like $\theta_1$ should be something like $\sqrt{m}\cdot \log^4 n/\eps^2$}
\item Repeat the following $t= \poly(\log n/\eps)$
% \Theta(\log^4 n/\eps^5)$ 
times: 
%\commentt{DR: verify }
  \begin{enumerate}
     \item Select a vertex $v$ uniformly at random and query its degree.  If $d(v) > \theta_0$, then continue to the next iteration.
     \item Perform a \emph{restricted} BFS starting from $v$ to depth $4$, %\commentt{the analysis seems to need 4} 
     where the BFS is restricted in the following sense.
         \begin{enumerate}
           \item Whenever a vertex $u$ is reached such that $d(u) \leq \theta_0$, all its neighbors are queried.
           \item Whenever a vertex $u$ is reached such that $d(u) > \theta_0$ and $u$ is reached from a vertex $u'$ such that $d(u') \leq \theta_0$, there are two sub-cases. If $d(u) \leq \theta_1$, then  all of $u$'s neighbors are queried. Otherwise,  $\theta_1$ neighbors of $u$ are selected uniformly at random. \label{step:C6_rand_nbr}
               %\commentt{Or should we use $\ell$ here?}
           \item Whenever a vertex $u$ is reached from a vertex $u'$ such that both $d(u)>\theta_0$ and $d(u') > \theta_0$, the BFS does not continue from $u$.
         \end{enumerate}
     \item If a $C_6$ is detected, then return `Reject' (and terminate).
%%%%%%%%%%%%%%%%%%%
\iffalse
     \item Call the procedure \textsf{One-high(v,n,\theta_0)}.
     \item Call the procedure \textsf{Two-high(v,n,\theta_0)}.
   %   \commentt{The case when there are two high-degree vertices on most $C_6$s.}
     \item Call the procedure \textsf{Three-high(v,n,\theta_0)}.
     \item If any one of the procedures detected a $C_6$, then return `Reject' (and terminate).
\fi
    \end{enumerate}
  \item Return `Accept'.
%%%%%%%%%%%%%%%%
\end{enumerate}

% \commentt{DR: initially I had calls to three procedures, depending on the number of high-degree vertices on the $C_6$. Maybe will want to return to that (?)}

%%%%%%%%%%%%%%%%%%%%%%%
\iffalse
\commentt{DR: say somewhere that if we encounter the same vertex $v\in L$ more than once, then we consider the same pair of incident edges. Alternatively, we can consider a version of the procedures where we always look at all neighbors of vertices in $L$. In such a case maybe we can actually give one algorithm which performs a kind of modified BFS to depth 3 from $v$: For vertices in $L$ (with degree  $\leq \theta_0$) we look at all their neighbors; for vertices not in $L$ that were reached from a vertex in $L$, if their degree is at most $\sqrt{n}$ (or so), then we look at all their neighbors, and otherwise we same $\sqrt{n}$ of their neighbors. (For vertices not in $L$ that were reached from a vertex not in $L$ we do not continue the search.)}

\smallskip\noindent
\textsf{One-high}($v,n,\theta_0$)
  \begin{enumerate}
     \item Select two random neighbors of $v$, denoted $u_1$ and $u_2$, respectively.
     \item If $d(u_1) \leq \theta_0$ and $d(u_2) \leq \theta_0$, then select a random neighbor $w_1$ of $u_1$ and a random neighbor $w_2$ of $u_2$.
      \item If $d(w_1) \leq \theta_0$ and $d(w_2) \leq \theta_0$, then select a random neighbor $x_1$ of $w_1$ and a random neighbor $x_2$ of $w_2$.
   \item If a $C_6$ is detected (since there is one path $v,u_1,w_1,x_1$ and another path $v,u_2,w_2,x_2$ where $x_1=x_2$, then return `Found', otherwise return `Not Found'.
  \end{enumerate}


%  Perform a restricted BFS starting from $v$ to depth 5, where the restriction is that the BFS does not continue from a vertex with degree greater than $\theta_0$. If a $C_6$ is detected, then return `Reject' (and terminate).

\smallskip\noindent
\textsf{Two-high}($v,n,\theta_0$)
  \begin{enumerate}
     \item Select two random neighbors of $v$, denoted $u_1$ and $u_2$, respectively.
  \end{enumerate}

\smallskip\noindent
\textsf{Three-high}($v,n,\theta_0$)
\begin{enumerate}
   \item Select two neighbors of $v$ uniformly at random, and denote them by $u_1$ and $u_2$ respectively.
   \item %Perform degree queries on $u_1$ and $u_2$.
          If $d(u_1)> \theta_0$ and $d(u_2) > \theta_0$, then do the following for each $b\in \{1,2\}$:
   \begin{enumerate}
     \item  Select $\min\{d(u_b),\sqrt{n}\}$ random neighbors of $u_b$.
     \item For each selected neighbor $w$ of $u_b$,  if $d(w) \leq \theta_0$, then select a random neighbor $x$ of $w$. \commentt{DR: to be completely consistent with the reduction, we need to check that $w$ selected $u_b$ as one of its neighbors, or maybe do this only in the analysis}
   \end{enumerate}
%   \item Obtain all neighbors of $v$ by performing $d(v)$ neighbor queries.
%   \item For each neighbor $u$ perform a degree query.
%   \item For every neighbor $u$ such that $d(u) > \theta_0$ do the following.
%       \begin{enumerate}
%            \item If $d(u) \leq \sqrt{n}$, then perform all neighbor queries from $u$, otherwise select, uniformly at random, $\sqrt{n}$ neighbors of $u$.
%            \item For each (selected) neighbor $w$ of $u$, query its degree. If $d(w) \leq \theta_0$ then obtain all %neighbors of $w$.
%       \end{enumerate}
   \item If a $C_6$ is detected (since there is one path $v,u_1,w_1,x$ and another path $v,u_2,w_2,x$ for some vertex $x$), then return `Found', otherwise return `Not Found'.
\end{enumerate}
\fi
%%%%%%%%%%%%%%%%%%%%%

%\begin{claim}
%\end{claim}

Since the algorithm rejects only when it detects a $C_6$, it always accepts when the graph $G$ is $C_6$-free, and hence, from this point on, we focus on the case that $G$ is $\eps$-far from $C_6$-free.
Let $\calC$ be a maximal set of edge-disjoint $C_6$s in $G$ where for every $C_6$ in $\calC$, there are no edges in which both endpoints have degree greater than $\theta_0$
(where $\theta_0 = 4\alpha/\eps$, as defined in the algorithm).
% so that $|\calC| > \eps m/6$ (where $m = \Theta(n)$ since $G$ has constant arboricity).\talya{shouldn't this be just $\eps m$?} \commentt{I don't think so: supposedly, in order to cover all $C_6$s (and not only the set of disjoint ones) we might need to remove all edges from this set}
% Let $\calC_{>\theta_0}$ be the subset of $C_6$s in $\calC$ that contain an edge between two vertices with degree greater than $\theta_0$ (where $\theta_0 = c_0\alpha/\eps$, as defined in the algorithm). 
Since the arboricity of $G$ is at most $\alpha$, 
% for $c_0=4$, 
we have that $|\calC| \geq \eps m/12$. 
% \commentt{changed here a bit, also in the same place for $C_5$}
% $|\calC'| < \eps m/c'$ %\commentt{DR: make more precise}.
% Observe that the algorithm only tries to detect $C_6$s in $\calC \setminus \calC'$ (as the BFS does not continue once an edge between two vertices with degree greater than $\theta_0$ is traversed).

Let $L$ denote the set of vertices with degree at most $\theta_0$, let $M$ denote the vertices with degree greater than $\theta_0$, and let $H$ denote those with degree greater than $\theta_1$ (where $L$ stands for \emph{low}, $M$ for \emph{medium}, and $H$ for \emph{high}).
% For each vertex $v \in L$, consider selecting two neighbors of $v$, uniformly at random. An edge between two vertices in $L$ is selected if both its endpoints selected each other. \commentt{DR: This seems too wasteful - maybe apply this only in the case of three high-degree vertices, and in the other two cases look at all neighbors of vertices in $L$.}
%  For each $C_6$ in $\calC \setminus \calC'$, the probability that all its edges were selected is at least $1/(\theta_0)^6$. By a multiplicative Chernoff bound, with probability $1-\exp(-\Omega(m))$, the number of such $C_6$s is at least $c''\eps^7 m$ for a constant $c''$. From this point on we condition on this event. While the algorithm selects the pair of neighbors of each
 We partition the $C_6$s in $\calC$ into three subsets:  $\calC_1$ contains those with at most one vertex in $M$,  and for $b\in \{2,3\}$, $\calC_b$ contains those with exactly $b$ vertices in $M$. For each $b\in \{1,2,3\} $, we consider the case that $|\calC_b| \geq \eps m/36$, where for at least one of the three this lower bound should hold. %\sum_{b'}|\calC_{b'}|/3$.

\begin{claim}\label{clm:C6-1}
If $|\calC_1| \geq \eps m/36$, then the algorithm will detect a $C_6$ with probability at least $2/3$.
\end{claim}

\begin{proof}
Consider the following set of \emph{useful} vertices $U$. For each $\rho\in \calC_1$, if all vertices on $\rho$ belong to $L$, then they all belong to $U$. If there is a vertex $u$ on $\rho$ that belongs to $M$, then the vertex $v\in L$ at distance $3$ from $u$ (on $\rho$), belongs to $U$ (that is, the vertex $v$  that is opposite to $u$ on $\rho$). Since $|\calC_1| \geq \eps m/36$, and each vertex in $L$ belongs to at most $\theta_0/2$ edge-disjoint $C_6$s in $\calC_1$, we have that $|U| = \Omega(\eps n)$. Therefore, with probability at least $2/3$, some vertex $v\in U$ is selected in at least one of the iterations of the algorithm. Conditioned on this event, the restricted BFS starting from $v$ will detect a $C_6$.
\end{proof}

\begin{claim}\label{clm:C6-2}
If $|\calC_2| \geq \eps m/36$, then the algorithm will detect a $C_6$ with probability at least $2/3$.
\end{claim}

\begin{proof}
Recall that $\calC_2$ is the set of $C_6$s in $\calC$ with exactly two vertices in $M$, i.e., that have degree greater than $\theta_0$. We consider several sub-cases depending on the number of  these vertices that are also in $H$. 
Let  $\calC_{2,1}$ consists of those $C_6$s in $\calC_2$ that contain at most one vertex in $H$, and $\calC_{2,2}$ those that contain two vertices in $H$, and note that $\calC_2=\calC_{2,1}\cupdot \calC_{2,2}$ so that at least one of these subsets is of size at least $|\calC_2|/2$. 

If $|\calC_{2,1}| \geq |\calC_2|/2$, then let $U$ denote those vertices in $L$ that reside on $C_6$s in $\calC_{2,1}$ that either do not have any vertex in $H$, or are at distance at least two from the single vertex in $H$ on the $C_6$. Similarly to the proof of Claim~\ref{clm:C6-1}, the algorithm selects a vertex from $U$ with probability at least $2/3$, and conditioned on this event, a $C_6$ is detected. 
% \commentt{in the future add drawings}.

We now turn to the case that $|\calC_{2,2}| \geq |\calC_2|/2$. We further partition the $C_6$s in $\calC_{2,2}$ according to the distance on the $C_6$ between the two vertices in $H$, which is either $2$ or $3$, and denote the two subsets by $\calC_{2,2,2}$ and $\calC_{2,2,3}$, respectively. In what follows we analyze the case that $|\calC_{2,2,2}| \geq |\calC_{2,2}|/2$.
The case that $|\calC_{2,2,3}| \geq |\calC_{2,2}|/2$, is analyzed similarly.


% Similarly to what we did in the upper bound for $C_4$, we run a coloring process, where here we only color vertices in $H$. Initially all vertices in $H$ and all edges on $C_6$s in $\calC_{2,2,2}$ are colored green. As long as there is a vertex $u\in H$ that is colored green such that the number of green edges incident to $u$ is less than $(\eps/600) d(u)$  then
% %commentt{Set $\eps_1$} 
% we color $u$ red, and we also color red all green edges on $C_6$s that $u$ resides on. Let $\calC'_{2,2,2} \subseteq \calC_{2,2,2}$ denote the resulting subset of $C_6$s in which both vertices that belong to $H$ remain colored green so that $|\calC'_{2,2,2}| \geq \eps m/300$, and let $H'$ denote the vertices in $H$ that remain colored green.

For each pair of vertices $u,u' \in H$, let $\calC(u,u')$ denote the subset of $C_6$s in $\calC_{2,2,2}$ on which $u$ and $u'$ reside. 
We shall say that the set $\calC(u,u')$ is \emph{useful} if 
$|\calC(u,u')| \geq \frac{(\eps/600)\cdot \max\{d(u),d(u')\}}{|H|}$, otherwise it is \emph{un-useful}.
By this definition, the number of  $C_6$s 
in $\calC_{2,2,2}$ that belong to un-useful sets is at most 
\begin{equation}
    \sum_{u\in H}\sum_{u'\in H:d(u')\leq d(u)}\frac{(\eps/600)\cdot d(u)}{|H|}\leq 
    \eps m/300 \;.
\end{equation}
Let $\calC'_{2,2,2}$ denote the subset of $C_6$s in $\calC_{2,2,2}$ that belong to useful sets, so that
$|\calC'_{2,2,2}| \geq \eps m/300$.
% Recall that on each $C_6$, the two vertices have a common neighbor in $L$. We % shall say that the set $\calC'(u,u')$ is \emph{useful} if the number of common neighbors of $u$ and $u'$ is at least  
%\commentt{DR: DOUBLE CHECK, in particular the issue of max. I think it is Ok, since this way the vertex among $u$ and $u'$ that ``pays'' for removing the un-useful set is the one with the higher degree}

The algorithm ``tries" to detect $C_6$s of this form by starting from a vertex $v$ that is at distance $2$ from both $u$ and $u'$ on some $C_6$ in $\calC(u,u')$, and 
obtaining a path of length 2 and a path of length 4 that collide on either $u$ or $u'$. Assume w.l.o.g. that $d(u')\leq d(u)$. 
Recall that except for $u$ and $u'$, all other vertices on the $C_6$s in $\calC(u,u')$ are in $L$. Therefore, a restricted BFS starting from $v$ will surely reach both $u$ and $u'$ at depth 2. Then, for the algorithm to reach $u$ from  $u'$ in two more steps, it suffices that one of their common neighbors in $L$ is selected in Step~\ref{step:C6_rand_nbr} among the $\theta_1$ sampled random neighbors of $u'$. In such a case the restricted BFS will surely reach $u$  via $u'$,   and thus obtain a $C_6$.
\iffalse
Consider any pair $u,u'\in H$ such that $\calC'(u,u')$ is useful, and a BFS
that starts from a vertex $v$ that is at distance $2$ from $u$ and $u'$  on some $C_6$ in $\calC'(u,u')$.
(so that the paths from $v$ to $u$ and $u'$ on this $C_6$ pass through vertices in $L$). 
The probability that a $C_6$ is detected when performing a restricted BFS from $v$, equals the probability that 
once $u$ and $u'$ are reached from $v$, one of their common neighbors in $L$ is selected among the $\theta_1$ 
random neighbors of $u'$ in Step~\ref{step:C6_rand_nbr}. 
\fi
Since the number of common neighbors of $u$ and $u'$ is at last $|\calC(u,u')|$,
the probability of this event is at least 
$$1-\left(1-\frac{(\eps/600)\cdot \max\{d(u),d(u')\}}{|H|\cdot \min\{d(u),d(u')\}}\right)^{\theta_1}\;.$$
Since $|H| \leq 2m/\theta_1$, by the setting of $\theta_1$ (for a sufficiently large constant $c_1$ and recalling that the graph has constant arboricity, so that $m=O(n)$), this probability is at least $9/10$. Combining this with the probability of the selection of some vertex $v$ in $L$ as defined above (which occurs with probability $\Omega(\eps^2)$), the claim follows. 
% \commentt{ more details, notation (I guess in the future..}
%%%%%%%%%
%%%%%%%%%%%%%%%
\iffalse
Finally we address the case that $|\calC_{2,2,3}| \geq |\calC_{2,2}|/2$. We apply a coloring process which is the same as the one applied in the (sub-)case that $|\calC_{2,2,2}| \geq |\calC_{2,2}|/2$, except that here we start by coloring green all edges on $C_6$'s in $\calC_{2,2,3}$, and let $\calC'_{2,2,3} \subseteq \calC_{2,2,3}$ denote the resulting subset of $C_6$s in which both vertices that belong to $H$ remain colored green.
\fi
%%%%%%%%%%%
\end{proof}

\begin{claim}\label{clm:C6-3}
If $|\calC_3| \geq \eps m/36$, then the algorithm will detect a $C_6$ with probability at least $2/3$.
\end{claim}


% (as well as a similar claim for testing $C_8$-freeness), 

\begin{proof}
%\commentt{5.7: I started making some changes here}
Recall that by the definition of $\calC_3$, for each $C_6$ in $\calC_3$ there are three vertices that belong to $M$, and three that belong to $L$, where each vertex in $M$ neighbors two vertices in $L$.
We start by selecting a subset of $\calC_3$, denoted $\calC_3'$, such that each vertex in $L$ belongs to at most a single $C_6$ in $\calC_3'$. Since each vertex in $L$ has degree at most  $\theta_0 = O(1/\eps)$, if we select $\calC_3'$ in a simple iterative greedy manner, then we get that
$|\calC_3'| \geq \eps_3 m$, where $\eps_3 = \Omega(\eps^2)$. Let $M'\subseteq M$ be the subset of vertices in $M$ that reside on $C_6$s in $\calC_3'$ and let $L'\subseteq L$ be the corresponding subset of vertices in  $L$. Let the edges on the
$C_6$s in $\calC_3'$ be colored green in $G$.

We next define an auxiliary multi-graph $G'$ 
% whose set of vertices is a subset of$M$, and 
    where (some of the) triangles in $G'$ correspond to $C_6$s in $G$. 
%
For each vertex in $M'$  we have a vertex in $G'$. In addition, there is a special \emph{ground} vertex, which we denote by $v_g$. For each pair of vertices, $u_1,u_2\in M'$ and for each length-2 green path $(u_1,v,u_2)$ (where $v\in L'$), there is an edge in $G'$. We think of this edge as being colored blue in $G'$ and labeled by the vertex $v$ in $L'$.
% and for each vertex $v\in L'$ such that there is a length-2 path between $u_1$ and $u_2$ passing through $v$ on some $C_6$ in $\calC_3'$, there is an edge between $u_1$ and $u_2$ in $G'$. We think of this as a \emph{blue} edge that is labeled by $v$.
%In order to define the set of edges in $G'$,
% We first select for each vertex $v\in L$ that resides on at last one $C_6$ in $\calC_3'$, a single pair of edges incident to $v$ on one of these $C_6$s, and color them blue. The choice of these pairs of edges is such that if we let $\calC'_3$ denote the subset of $C_6$s in $\calC_3$ for which all edges are blue, then $|\calC_3'| \geq \eps_3 m$, where $\eps_3 = \Omega(\eps^4)$. A simple probabilistic argument (based on randomly selecting each such pair of edges) shows that there is such a choice of blue edges. 
% \commentt{DR: give more details, in particular set $\eps_3$, which I think is of the order of $\eps^4$}
% We put an edge in $G'$ between two vertices $u$ and $u'$ in $M$ (that reside on some $C_6$ in $\calC_3$) for every length-2 blue path between them. 
In addition, for each  $u\in M'$, and each non-green incident edge that it has in $G$, there is an (``uncolored'') edge between $u$ and $v_g$.
By this construction, there are $|\calC_3'|$ blue edge-disjoint triangles in $G'$ (and the degree of every vertex $u\in M'$ in $G'$ is the same as its degree in $G$).

%%%%%%%%%%%%%%%%%%%%%%%%%%%%%%%%%%%
\iffalse
Here too we run a (``green-red'') coloring procedure, where initially all vertices in $M$ and all edges in $G'$ between them are colored green. As long as there is a vertex $v$ that is colored green and the number of green edges incident to it is smaller than $\eps_3 d(v)/4$ (where $d(v)$ is its degree in $G$), we select such a vertex, and color it and its incident green edges by red. At the end of this process,  $G'$ contains at least $|\calC_3'| - 2\eps_3 m \geq \eps_3 m/2$ edge-disjoint triangles in which  all vertices and edges are green.
\commentt{maybe we don't need to  perform this coloring here since we perform a related process inside the proof }
\fi
%%%%%%%%%%%%%%%%%

% \commentt{DR: we can't use the claim exactly as is for $C_8$. Will think about this later. It probably should be about a general multigraph $G'$, but we also have the degree of the vertices in $G$.}
\begin{claim}\label{clm:tri-G-prime}
Let $G'$ be a multigraph over $m' \leq m$ edges, with $m'_b$ blue edges, and 
at least $\eps_3 m$ edge-disjoint blue triangles. 
Let $D$ be any distribution over edges in $G'$ such that for each blue edge in $G'$, the probability that it is selected according to $D$ is %at least $\eps'/m'$ 
$\Omega(1/m'_b)$.
% where $m'$ is the number of blue edges in $G'$. 
Consider the following process.
First we select a blue edge $(u,u')$ according to $D$. 
Next, for each of $u$ and $u'$, if its degree %(in $G'$ as well as in $G$) 
is at most $\theta_1$, then we query all its neighbors in $G'$, and otherwise
we select $\theta_1$  random neighbors in $G'$.
 The probability that the process obtains a triangle $(u,u',u'')$ in $G'$ where $u''\in M'$ is 
%$\Omega(\eps'\cdot \eps_3/\log^4 n)$.
$\Omega(\eps_3/\log^4 n)$.
\end{claim}

\begin{proof}
Let $\calT$ denote a maximal set of edge-disjoint triangles in $G'$ whose edges are all colored blue, so that $|\calT| \geq \eps_3 m$. Let $\calT_1$ denote the subset of triangles in $\calT$ that contain at most one vertex in $H$, and let $\calT_2$ denote those that contain at least two vertices in $H$. If $|\calT_1| \geq |\calT|/2$, then we are done:  conditioned on selecting an edge between two vertices in $M'\setminus H$ that belongs to a triangle in $\calT_1$, the procedure will detect this triangle, as it queries all neighbors of the selected edge's endpoints, and the probability of selecting such an edge is %$\Omega(\eps'\cdot \eps_3)$.
$\Omega(\eps_3)$.

From this point on we address the case that $|\calT_2| \geq |\calT|/2$. 
For a vertex $u\in M'$ we use $d(u)$ to denote its degree in $G'$. %\emph{in $G$}.
% \commentt{DR: need to give intuition at some point to why we do what we do}.
For each pair of vertices $u_1,u_2\in M'$, let $E(u_1,u_2)$ denote the set of edges between $u_1$ and $u_2$ that reside on triangles in $\calT_2$ (so that in particular, they are all blue), and let  $e(u_1,u_2) = |E(u_1,u_2)|$. 
% \commentt{maybe do without $e(\cdot,\cdot)$ notation and just use $|E(\cdot,\cdot)|$? Note that the text below (until the bucketing) is not really necessary technically, but may help explain where we are heading.}
Suppose the procedure described in the claim selects an edge $(a,b)\in E(u_1,u_2)$ where $a,b\in H$. Recall that in such a case, $\theta_1$ random neighbors of $a$ and $\theta_1$ random neighbors of $b$ are sampled.
For $i \in [\theta_1]$ and $j\in [\theta_1]$, let $\chi_{i,j}(a,b)$ be a Bernoulli random variable that is $1$ if and only if the $i$th randomly selected neighbor of $a$ equals the $j$th randomly selected neighbor of $b$, where this neighbor is not $v_g$ (thus detecting a blue triangle).
Note that by this definition, the expected number of triangles detected (conditioned on  selecting $(a,b)$) is
$\sum_{i,j\in [\theta_1]}\chi_{i,j}(a,b)$,
%Then the probability that a triangle is detected is when selecting a random neighbor of
where
\begin{equation}\label{eq:a-b-c-1}
    \Pr[\chi_{i,j}(a,b)=1]= \sum_{c\in M}\frac{e(a,c)}{d(a)}\cdot \frac{e(b,c)}{d(b)}\;.
\end{equation}
% DR: we had $c\neq a,b$ but this is not necessary, since if $c=a$, then $e(a,c)=0$ and analogously for $b$
A simple but important observation is that, since each of the $e(a,b)$ edges in $E(a,b)$ participates in some triangle $(a,b,c)$ in $\calT_2$ where these triangles are edge-disjoint, $e(a,b)$ is upper bounded by
% the sum over all $c$ for which $e(a,c)$ and $e(b,c)$ are both non-zero of $\max\{e(a,c)\cdot e(b,c)\}$
$\sum_{c \in M} \min\{e(a,c),e(b,c)\}$, which in turn is upper bounded by
$\sum_{c\in M} e(a,c)\cdot e(b,c)$. 
That is, 
$$\sum_{c\in M} e(a,c)\cdot e(b,c)\geq \sum_{c \in M} \min\{e(a,c),e(b,c)\} \geq e(a,b)\;.$$
Therefore,
\begin{equation}\label{eq:a-b-c-2}
    \Pr[\chi_{i,j}(a,b)=1] \geq \sum_{c\in M}\frac{e(a,b)}{d(a)\cdot d(b)}\;.
\end{equation}
We next show that there exists a relatively large set of edges between  pairs $(a,b)$ where $a,b \in H$, for which
% the expression on the right-hand-side of Equation~\eqref{eq:a-b-c-2} 
 $\Pr[\chi_{i,j}(a,b)=1] = \tilde{\Omega}(1/n)$. We shall refer to such pairs $a,b$  as \emph{useful} pairs. %\commentt{be more precise}
This will imply  that the expected number of triangles detected when selecting an edge between one of these pairs is $\tilde{\Omega}(\theta_1^2/n)$. %  \talya{I forgot why we get $\frac{\theta_1^2}{n}?$ from the above I can only see that we get $\frac{e(a,b)}{d(a)d(b)\geq \frac{\theta_1^2 }{n^3}}$. } 
% \commentt{it is because there are $\theta_1^2$ such pairs}
By the setting of $\theta_1$, this is greater than 1. %\talya{aren't we talking about a probability?} 
%\commentt{No, we wrote expected number}
Our definition of  this set of useful pairs  will also ensure that the variance is bounded, allowing us to deduce that a triangle is detected with sufficiently high probability. Details follow.


We start by partitioning the set $\calT_2$ into subsets as follows. First, for each triangle $\Delta \in \calT_2$, we 
denote it by $(a_\Delta, b_\Delta, c_\Delta)$, where, without loss of generality, $a_\Delta, b_\Delta\in H$. For every four integers $x_{\tiny A},x_{\tiny B},y_{\tiny A,C},y_{\tiny B,C}$,  the subset $\calT_2(x_{\tiny A},x_{\tiny B},y_{\tiny A,C},y_{\tiny B,C})$ consists of all triangles $\Delta \in \calT_2$ for which the following holds.
$d(a_\Delta) \in [2^{x_{\tiny A}-1}\cdot \theta_1 , 2^{x_{\tiny A}}\cdot \theta_1)$, $d(b_\Delta) \in [2^{x_{\tiny B}-1}\cdot \theta_1,  2^{x_{\tiny B}}\cdot \theta_1)$, $e(a_\Delta,c_\Delta) \in [2^{y_{\tiny A,C}-1}, 2^{y_{\tiny A,C}})$, and $e(b_\Delta,c_\Delta) \in [ 2^{y_{\tiny B,C}-1},  2^{y_{\tiny B,C}})$.
% $2^{x_{\tiny A}-1}\cdot \theta_1 \leq d(a_\Delta) < 2^{x_{\tiny A}}\cdot \theta_1$, 
% $2^{x_{\tiny B}-1}\cdot \theta_1 \leq d(b_\Delta) < 2^{x_{\tiny B}}\cdot \theta_1$, 
% $2^{y_{\tiny A,C}-1}\leq e(a_\Delta,c_\Delta) < 2^{y_{\tiny A,C}}$, and
% $2^{y_{\tiny B,C}-1}\leq e(b_\Delta,c_\Delta) < 2^{y_{\tiny B,C}}$. 
Let $\calT^*_2 = \calT_2(x^*_{\tiny A},x^*_{\tiny B},y^*_{\tiny A,C},y^*_{\tiny B,C})$ be the largest subset, so that $|\calT^*_2| \geq \eps_3 m/\log^4 n$. 
Let $A^* \eqdef \{a_\Delta: \Delta \in \calT^*_2\}$ and $B^* \eqdef \{b_\Delta:\Delta \in \calT^*_2\}$. 
%\commentt{say something about the possibility that these sets are not disjoint?}
For each pair of vertices $a\in A^*$ and $b\in B^*$, let $E^*(a,b)$ be the subset of
edges between $a$ and $b$ such that there exists $\Delta\in \calT^*_2$ for which
$a = a_\Delta$ and $b = b_\Delta$, let $e^*(a,b) = |E^*(a,b)|$
and let $C^*_{a,b} \eqdef \{c_\Delta\,:\, a = a_\Delta, b=b_\Delta, \Delta \in \calT^*_2\}$. 
% \commentt{I didn't find the need to define $A$ and $B$, but maybe helpful for understanding?}
We have:
\begin{equation}\label{eq:sum-e-star}
s^* \eqdef \sum_{a\in A^*, b\in B^*} e^*(a,b) \geq \frac{\eps_3 m}{\log^4 n}\;,
\end{equation}
and for every $a \in A^*$ and $b\in B^*$, similarly to what was argued before regarding $e(a,b)$,
\begin{equation}\label{eq:e-star-ub}
\sum_{c\in C^*(a,b)} e(a,c)\cdot e(b,c) \geq e^*(a,b)\;.
\end{equation}
%%%%
\iffalse
For each $a \in A^*$ let $d^*(a) \eqdef |\{\Delta \in \calT^*_2 : a = a_\Delta\}|$ and
for $b \in B^*$, define $d^*(b)$ similarly. \commentt{this is a bit different from the notes -- double-check} 
Let $\calT^{**}_2$ consist of those triangles $\Delta$ in $\calT^*_2$ for which both 
$d^*(a_\Delta) \geq (\eps_3/(XX\log^4 n))\cdot d(a_\Delta)$ and 
$d^*(b_\Delta) \geq (\eps_3/(XX\log^4 n))\cdot d(b_\Delta)$, so that
$|\calT^{**}_2| \geq  |\calT^*_2| - 4 (\eps_3/(XX\log^4 n))\cdot m = \Omega(\eps_3 m/\log^4 n)$.
Let $A^{**} \eqdef \{a_\Delta: \Delta \in \calT^**_2\}$ and $B^{**} \eqdef \{b_\Delta: \Delta \in \calT^**_2\}$,
so that 
\begin{equation}\label{eq:sum-e-double-star}
s^{*} \eqdef \sum_{a\in A^{**}, b\in B^{**}} e^*(a,b) = \Omega(\eps_3 m/\log^4 n)\;.
\end{equation}
\fi
%%%%%%%%%%%%%%%%
Let 
\begin{equation}\label{eq:def-P-star}
P^{*}\eqdef \left\{(a,b)\,:\, a\in A^{*},b\in B^{*}, e^*(a,b) \geq \frac{s^*}{2\cdot |A^{*}|\cdot|B^{*}|} \right\}\;.
\end{equation}
Since 
\begin{equation}\sum_{(a,b)\in A^*\times B^* \setminus P^{*}} e^*(a,b) \leq |A^*|\cdot |B^*| \cdot \frac{s^*}{2\cdot |A^{*}|\cdot|B^{*}|} = s^*/2\;,
\end{equation}
by Equation~\eqref{eq:sum-e-star},
\begin{equation}\label{eq:sum-P-star}
\sum_{(a,b)\in P^{*}} e^*(a,b) \geq \frac{\eps_3 m}{2\log^4 n}\;.
\end{equation}

By Equation~\eqref{eq:sum-P-star} and the premise of the claim regarding $D$, if we select an edge according to $D$, then the probability that we obtain an edge in $E^*(a,b)$ for $(a,b)\in P^{*}$ is %$\Omega(\eps'\eps_3/\log^4 n)$.
$\Omega(\eps_3/\log^4 n)$.
Fixing any such pair $(a,b)$, we slightly modify the definition of the random variables $\chi_{i,j}(a,b)$ so that $\chi_{i,j}(a,b)=1$ if and only if the $i$th random sampled neighbor of $a$ and the $j$th random sampled neighbor of $b$ are both a common neighbor $c \in C^*(a,b)$ (the original definition allowed any common neighbor $c$ (other than $v_g$)). %\talya{slightly changed wording}
We next upper bound the probability that $\sum_{i,j\in [\theta_1]}\chi_{i,j}(a,b) =0$. 
Let 
\begin{equation}\label{eq:mu-def}
\mu(a,b)\eqdef \sum_{c\in C^*(a,b)}\frac{e(a,c)}{d(a)}\cdot \frac{e(b,c)}{d(b)}\;,
\end{equation}
so that $\Ex[\chi_{i,j}(a,b)] = \Pr[\chi_{i,j}(a,b)=1] = \mu(a,b)$.
Since $(a,b)\in P^*$,
by Equations~\eqref{eq:sum-e-star}--\eqref{eq:def-P-star},
%Equation~\eqref{eq:e-star-ub}, , and 
%Equation~\eqref{eq:def-P-star},
% (where $P^{*}$ is as defined in Equation~\eqref{eq:def-P-star}),
\begin{equation}\label{eq:mu-lb}
\mu(a,b) \geq \frac{e^*(a,b)}{d(a)\cdot d(b)} \geq \frac{\eps_3 m}{2\log^4 n \cdot d(a)\cdot d(b)\cdot |A^*| \cdot |B^*|} \geq \frac{\eps_3 }{8 \log^4 n \cdot m}  \;,
\end{equation}
where we have used the fact that $d(a)\cdot |A^{*}| \leq 2m' \leq 2m$ and $d(b)\cdot |B^{*}| \leq 2m' \leq 2m$. 
% \talya{did we write somewhere that $m(G')=O(m)$?}
%\commentt{I guess we didn't explicitly. We only said that the degree of each vertex of $M$ in $G'$ equals its degree in $G$. Can add}
From this point on we shall use the shorthand $\chi_{i,j}$ for $\chi_{i,j}(a,b)$ and $\mu$ for $\mu(a,b)$.

At this point we would have liked to apply Corollary~\ref{cor:cheb1} to upper bound the probability that
$\sum_{i,j\in [\theta_1]} \chi_{i,j}=0$, but it does not exactly meet our needs, and hence we give a direct proof.
% Let $\bar{\chi}_{i,j} = \chi_{i,j} - \mu$, so that $\Ex[\bar{\chi}_{i,j}] = 0$.
First we observe that 
\begin{equation}
  \Ex\left[\sum_{i,j\in [\theta_1]}\chi_{i,j}\right] = \theta_1^2\cdot \mu\;.  
\end{equation}
By applying Claim~\ref{clm:var1} (see Appendix~\ref{app:prob}), with $s_1=s_2=\theta_1$ and using the notation $\mu_{1,2} = \Pr[\chi_{i,j_1}=\chi_{i,j_2}=1]$ (for every $i,j_1,j_1\in [\theta_1]$) and
$\mu_{2,1} = \Pr[\chi_{i_1,j} = \chi_{i_2,j} = 1]$ (for every $i_1,i_2,j\in [\theta_1]$) from that claim, we have that
\ForFuture{T: add elaboration}
\begin{equation}\label{eq:var-sum-chis}
    \Var\left[\sum_{i,j\in [\theta_1]}\chi_{i,j}\right] \leq \theta_1^2\cdot \mu + \theta_1^3 \cdot \mu_{1,2} + \theta_1^3 \cdot \mu_{2,1} + {\theta_1\choose 2}^2\cdot \mu^2 - (\theta_1^2\cdot \mu)^2\;.
\end{equation}
By Chebishev's inequality and the above two equations,
\begin{equation}\label{eq:C6-cheb2}
\Pr\left[\sum_{i,j\in [\theta_1]}\chi_{i,j}=0\right] \leq \frac{\Var\left[\sum_{i,j\in [\theta_1]}\chi_{i,j}\right]}{\left(\Ex\left[\sum_{i,j\in [\theta_1]}\chi_{i,j}\right]\right)^2}
\leq  \frac{1}{\theta_1^2\mu} + \frac{\mu_{1,2}}{\theta_1\cdot \mu^2} + \frac{\mu_{2,1}}{\theta_1\cdot \mu^2} \;.
\end{equation}
We next upper bound each of the three terms on the right-hand-side of Equation~\eqref{eq:C6-cheb2}.

By Equation~\eqref{eq:mu-lb} and the setting of $\theta_1=c_1\cdot n^{1/2}\log^2 n/\eps^2$, the first term contributes at most $1/6$ (for a sufficiently large constant $c_1$). %\commentt{DR: double check after setting $\theta_1$}


Turning to the second term, 
\begin{equation}
\mu_{1,2} = 
\Pr[\chi_{i,j_1}=\chi_{i,j_2}=1] 
  = \sum_{c\in C^*(a,b)} \frac{e(a,c)\cdot (e(b,c))^2}{d(a)\cdot (d(b))^2}\;.
\end{equation}
By the definition of $\mu$ in Equation~\eqref{eq:mu-def}
and again using Equations~\eqref{eq:sum-e-star}--\eqref{eq:def-P-star}, the second term is upper bounded by
\begin{eqnarray}
\frac{1}{\theta_1}\cdot \frac{d(a)\cdot \sum_{c\in C^*(a,b)} e(a,c)\cdot (e(b,c))^2 }
                           {\left( \sum_{c\in C^*(a,b)} e(a,c)\cdot e(b,c)\right)^2}  &\leq&
\frac{1}{\theta_1}\cdot \frac{d(a) \cdot |C^*(a,b)| \cdot 2^{y^*_{\tiny A,C}}\cdot 2^{2y^*_{\tiny B,C}}} 
                           {|C^*(a,b)|^2 \cdot 2^{2(y^*_{\tiny A,C}-1)}\cdot 2^{2(y^*_{\tiny B,C}-1)}} \\
 &=&      
 \frac{1}{\theta_1}\cdot \frac{ 16d(a) }
                           {|C^*(a,b)| \cdot 2^{y^*_{\tiny A,C}}} \\   
 &\leq&  \frac{1}{\theta_1}\cdot \frac{16 d(a) }{e^*(a,b)}       \\
 &\leq&    \frac{1}{\theta_1}\cdot \frac{8 \log^4 n \cdot d(a) \cdot |A^{*}|\cdot |B^{*}|}{\eps_3 m }   \\
 &\leq& \frac{16 \log^4 n \cdot |B^{*}|}{\eps_3 \theta_1}  \;.
\end{eqnarray}
Since $|B^*| \leq 2m/\theta_1$ (as $B^* \subseteq H$), by the setting of $\theta_1$, the contribution of this term is at most $1/6$ as well.

For the third term we have that 
\begin{equation}\label{eq:cheb-C6-case3}
\mu_{2,1} = 
  \Pr[\chi_{i_1,j_1}=\chi_{i_1,j_2}] = 
   \sum_{c\in C^*(a,b)} \frac{(e(a,c))^2\cdot e(b,c)}{d(a)\cdot d(b)}\;.
\end{equation}
By essentially the same argument as the one bounding the second term (replacing the roles of $a$ and $b$ ($A^*$ and $B^*$), this term contributes at most $1/6$ as well. 

We have thus shown that, conditioned on selecting an edge between a pair $(a,b) \in P^*$, the probability that we obtain a triangle is at least $1/2$, and Claim~\ref{clm:tri-G-prime} follows.
\end{proof}

\medskip
In order to complete the proof of
Claim~\ref{clm:C6-3} by applying Claim~\ref{clm:tri-G-prime}, we make the following observations. First, by the construction of $G'$, each vertex in $L'$  corresponds to exactly one blue edge in $G'$. In particular, if  we select a vertex in $L'$ and query all its neighbors (as done in the algorithm), then we effectively obtain an edge in $G'$. Since $|\calC_3'|= \Omega(\eps^2 m)$, we have that $|L'| = \Omega(\eps^2 n)$.
Therefore, the probability in each iteration that the algorithm selects a vertex in $L'$, and hence an edge in $G'$, is $\Omega(\eps^2)$. 
% (and the queries performed on its endpoints
% then the probability that we obtain a pair of vertices in $M$ that have an edge between them in $G'$ is at least $\eps'/m'$, for $\eps' = \Omega()$.

Next, consider each of the endpoints $u$ of such an edge. If $d(u)\leq \theta_1$, then our algorithm queries all the neighbors of $u$, and for each neighbor $v$ of $u$ such that $d(v) \leq \theta_0$, the algorithm queries all the neighbors of $v$. This implies that  the algorithm reaches all of the  neighbors $u'\in M'$ of $u$ in $G'$ (as in the corresponding case described in Claim~\ref{clm:tri-G-prime}). Similarly, if $d(u)> \theta_1$, then our algorithms samples $\theta_1$ random neighbors of $u$, and for each sampled neighbor $v$ such that $d(v) \leq \theta_0$, it queries all the neighbors of $v$. This implies that for every $u'\in M'$, the probability that it is reached from $u$ by such a length-2 path in $G$, is at least as large as the probability that it is reached in one random step from $u$ in $G'$.
Since blue triangles in $G'$ correspond to $C_6$s in $G$, Claim~\ref{clm:C6-3} follows.
\end{proof}

\medskip
The proof of Theorem~\ref{thm:C6-ub} follows directly from Claims~\ref{clm:C6-1}--\ref{clm:C6-3} (together with the description of the algorithm for obtaining the upper bound on the complexity of the algorithm).


\iffalse
\commentt{DR: or maybe we should prove a ``general reduction theorem'' and derive the upper bound for C6 and C8 from it, after proving results $C_3$ and $C_4$ freeness results in multigraphs}

\begin{theorem}\label{thm:gen-red-even}
Suppose that for a given integer $k\geq 3$, for every $3 \leq k' \leq k$, there is a testing algorithm for $C_{k'}$ freeness in general multigraphs that is given access to random edge queries and random neighbor queries \commentt{DR: be more precise here} whose query complexity is $g_{k'}(m',\eps)$ where $m'$ is the number of edges in the multigraph. Then there exists a testing algorithm for $C_{2k}$-freeness in graphs with constant arboricity whose query complexity is $O(\max_{k'\leq k} g_{k'}(n,\eps))$
\end{theorem} 
\fi


